\subsubsection{Sites cristallographiques interstitiels : \og critère stérique \fg{}}
%Marruco, Exercice 2.1

\setcounter{equation}{0}
\setcounter{Q}{0}

Cet exercice cherche à définir un critère stérique simple permettant de rendre compte de l'occupation d'un site interstitiels cubique, octaédrique et tétraédrique.

Dans tout l'exercice, on considérera que les cations (resp. anions) de rayon $R_\oplus$ (resp. $R_\ominus$) sont des sphères dures tangentes entre elles.

\paragraph*{Site cubique}

Dans la structure cubique simple (CS) est représentée dans un cubique d'arête \textbf{a} dont les sommets sont occupés par l'anion et le centre par le cation. 

\Q Représenter la maille légendée.
\solution{Voici une maille type \ce{CsCl}
\begin{center}
\includegraphics[width=0.48\textwidth]{Ex13/Fig01}
\end{center}
Le cation (resp. l'anion) est représenté en vert (resp. en jaune).
}

\Q Montrer que pour les sites cubiques : 
\begin{align}
\rho_C = \frac{R_\oplus}{R_\ominus} = \sqrt{3} - 1 \approx 0.732
\label{eq:13-01}
\end{align}
\solution{L'atome central occupe un site cubique de la maille cubique simple (primitive). On écrit deux relations : 
\begin{align*}
\text{Condition de contact} & \quad \sqrt{3} a = 2 R_\oplus + 2 R_\ominus \\
\text{Définition de }a & \quad 2 R_\ominus = a
\end{align*}
Ce qui permet de remonter à la relation de l'Eq. \ref{eq:13-01}.}

\paragraph*{Site octaédrique}

Dans la structure cubique face centrée (CFC), un des sites octaédriques se situe au centre d'un cube. Les sommets et le milieu des faces de celui-ci sont occupés par des cations, tandis que les anions sont situés au milieu des arêtes de longueur \textbf{a}.

\Q Représenter la maille légendée en y faisant figurer un site O.
\solution{Voici une maille type \ce{NaCl} 
\begin{center}
\includegraphics[width=0.48\textwidth]{Ex13/Fig02}
\end{center}
Le site au centre est un site octaédrique.\\
Bien entendu, dans une maille \ce{NaCl}, on peut aussi représenter l'anion dans une structure CFC avec les cations occupant les sites octaédriques. C'est plutôt cette représentation que l'on utilisera pour la question suivante (si on cherche la taille d'un site cationique).
}

\Q Montrer que pour les sites octaédriques : 
\begin{align}
\rho_O = \frac{R_\oplus}{R_\ominus} = \sqrt{2} - 1 \approx 0.414
\label{eq:13-02}
\end{align}
\solution{On écrit deux relations : 
\begin{align*}
\text{Condition de contact} & \quad a = 2 R_\oplus + 2 R_\ominus \\
\text{Définition de }a & \quad 4 R_\ominus = \sqrt{2} a
\end{align*}
Ce qui permet de remonter à la relation de l'Eq. \ref{eq:13-02}.}

\paragraph*{Site tétraédrique}
Les anions sont aux sommets d'un cube d'arête a/2 et répartis de manière à former un tétraèdre. Ils sont tangents le long de la diagonale d'une face. Le site tétraédrique est au centre du cube, au milieu du tétraèdre : 

\Q Représenter la maille légendée en y faisant figurer un site T.
\solution{Voici une maille type \ce{NaCl} 
\begin{center}
\includegraphics[width=0.48\textwidth]{Ex13/Fig03}
\end{center}
Le site indiqué par une étoile est un site tétraédrique.
}

\Q Montrer que pour les sites tétraédriques : 
\begin{align}
\rho_T = \frac{R_\oplus}{R_\ominus} = \sqrt{\frac{3}{2}} - 1 \approx 0.225
\label{eq:13-03}
\end{align}
\solution{On écrit deux relations : 
\begin{align*}
\text{Condition de contact} & \quad \frac{\sqrt{3}}{2} a = 2 R_\oplus + 2 R_\ominus \\
\text{Définition de }a & \quad 4 R_\ominus = \sqrt{2} a
\end{align*}
Ce qui permet de remonter à la relation de l'Eq. \ref{eq:13-03}.
}

\paragraph*{Bilan}
\Q En déduire une règle, que nous appellerons \og critère stérique \fg{} (dans le cadre de ce cours) permettant de définir l'occupation par un cation d'un site interstitiel.
\solution{On retrouve la première règle de Pauling qui définit la géométrie et la coordination d'un cation métallique en fonction du rapport $R_\oplus / R_\ominus$.}
\begin{comment}
\begin{center}
\begin{tabular}{ccc}
\toprule
Coordinence & Géométrie & Rapport \\
\midrule
[4] & tétraèdrique & $< 0.225$ \\
[6] & octaédrique & $0.225-0.414$ \\
[8] & cubique & $0.414-0.732$ \\
\bottomrule
\end{tabular}\end{center}}
\end{comment}
% \rho C > O > T 
 
